% This paragraph describes the actual current fine-grained table runs.
% They are GAM-EWCLoRA, not verified implementations of the manuscript's
% squared-gradient FS-EWC with trajectory-averaged retained anchors.
\paragraph{EWC and GAM-EWC settings.}
The EWC regularization coefficient $\lambda$ is benchmark-specific
and remains fixed throughout each task stream.
For Aircraft, Cars, CUB-200, Flowers, and 5Datasets, respectively,
EWCLoRA uses $\lambda=200$, $80{,}000$, $200{,}000$, $10$, and
$50{,}000$, whereas GAM-EWCLoRA uses $\lambda=2{,}000$ on all five
benchmarks. The EWCLoRA configurations correspond to the eligible
archived runs with the highest final average accuracy for each dataset,
selected retrospectively using final test performance.
Both methods use importance decay $\gamma=1$, at most 100
current-task mini-batches for task-end Fisher estimation, and
per-task trace normalization with numerical stabilizer $10^{-12}$.
EWCLoRA uses zero Fisher damping.

For GAM-EWCLoRA, the recorded triples
\texttt{(gam\_grad\_rho, gam\_grad\_norm\_rho, gam\_grad\_gamma)}
are $(0.01,0.10,0.02)$ on Aircraft,
$(0.05,0.20,0.03)$ on Cars, and $(0.20,0.20,0.10)$ on both
CUB-200 and 5Datasets.
These are the gradient-probe displacement, gradient-difference ascent
displacement, and GAM gradient-decomposition coefficient, respectively;
the last coefficient is distinct from Fisher damping and importance decay.
The corresponding GAM coefficients $\beta_1,\beta_2,\beta_3$ are all 1,
adaptive perturbations are disabled, and the flatness-direction mixing
coefficient is 1. The Flowers archive does not explicitly preserve the
GAM probe and decomposition settings, so they are reported as unrecorded.
Task-end Fisher estimation uses the final trained checkpoint, and the
resulting importance is consolidated with the inherited statistic for
the next task. These archived runs do not perform trajectory checkpoint
averaging before Fisher estimation.
